%%%%%%%%%%%%%%%%%%%%%%%%%%%%%%%%%%%%%%%%%%%%%%%%%%
% Compiler: pdfLaTeX
% TeX Live version: 2024
%%%%%%%%%%%%%%%%%%%%%%%%%%%%%%%%%%%%%%%%%%%%%%%%%%

% \ifdefined\ifhandout\else
\newif\ifhandout
% \fi

% NOTE: for the final version, set \handouttrue
% \handoutfalse
\handouttrue

%%%%%%%%%%%%%%%%%%%%%%%%%%%%%%%%%%%%%%%%%%%%%%%%%%
% unofficial phase abstract
%%%%%%%%%%%%%%%%%%%%%%%%%%%%%%%%%%%%%%%%%%%%%%%%%%

% % -*- Mode:TeX -*-

\documentclass{cinc-abstract}

\usepackage{graphicx}
\usepackage{xcolor}
\usepackage{relsize}
\usepackage{pifont}
\usepackage{currfile}

\newcommand\wordcount{\input{|"texcount -inc -sum -0 -utf8 -ch -template={SUM} \currfilepath"}}


\begin{document}

% The title is set in 14-point Helvetica bold
\title{To write}

% The rest of the title block is set in 12 points Helvetica
\author {Hao Wen, Jingsu Kang\\ % First name, initials and surnames, no ``and''
\ \\ % leave an empty line between authors and affiliation
China Agricultural University\\  % gives affiliation of the first author only
Beijing, China} % city, [state or province,] country only

\maketitle

%%%%%%%%%%%%%%%%%%%%%%%%%%%%%%%%%%%%%%%%%%%%%%%%%%%%%%%%%%
% NOTE: The body of the abstract (exclusive of the title, authors, and authors’ affiliations) can be up to 300 words at most
%%%%%%%%%%%%%%%%%%%%%%%%%%%%%%%%%%%%%%%%%%%%%%%%%%%%%%%%%%


Aim: to write

Methods: to write

Results: to write

Conclusion: to write


% \ifhandout
% % do nothing
% \else
% \begin{center}
% {\larger[2]\color{red} \ding{43}\ding{43}  Totally \wordcount words \reflectbox{\ding{43}\ding{43}}}
% \end{center}
% \fi


\end{document}



%%%%%%%%%%%%%%%%%%%%%%%%%%%%%%%%%%%%%%%%%%%%%%%%%%
% conference short paper
%%%%%%%%%%%%%%%%%%%%%%%%%%%%%%%%%%%%%%%%%%%%%%%%%%

%%%%%%%%%%%%%%%%%%%%%%%%%%%%%%%%%%%%%%%%%%%%%%%%%%%%%%%%%%%%%%%%%%%%%%%%%%%%%%%%

% preamble

\documentclass[twocolumn]{cinc}

\usepackage[T1]{fontenc}
\usepackage{mathtools, graphicx, xurl, hyperref}
% \usepackage{tikz-cd}
\usepackage{tikz}
\usetikzlibrary{shapes, arrows.meta, decorations.pathmorphing, decorations.pathreplacing, backgrounds, positioning, fit, calc, shadows, shapes.misc}
% \tikzexternalize[prefix=tikz/,optimize command away=\includepdf]
\usepackage{boldline, multirow}
\usepackage{booktabs} % Allows the use of \toprule, \midrule and \bottomrule in tables
\usepackage{colortbl}
\usepackage{subcaption}
\usepackage{relsize}
\usepackage{fontawesome5}
\usepackage[ruled, vlined]{algorithm2e}

\usepackage{tabularx} % for 'X' col. type and 'tabularx' environment
\usepackage{ragged2e} % for '\RaggedRight' macro
\newcolumntype{L}{>{\RaggedRight}X} % disable full justification

\usepackage{adjustbox}

% \tikzstyle{smallcircleblock} = [circle, draw, text width = 0.7em, text centered]

% \tikzstyle{wideblock} = [rectangle, draw, text width = 6.5em, text centered, rounded corners, inner sep = 7pt, minimum height = 1.0em]

\newcommand\ImageNode[4][]{
  \node[#1] (#2) {\includegraphics[#3]{#4}};
}
\definecolor{arrowblue}{RGB}{98,145,224}


\newcolumntype{g}{>{\columncolor[gray]{0.8}}l}


\newif\ifcoloredtext
\coloredtextfalse

\newif\ifboxednn
\boxednntrue


% --- unify spacing between floats and text (single-/double-column) ---
\setlength{\textfloatsep}{6pt plus 2pt minus 2pt}   % text <-> single-column floats (table/figure)
\setlength{\intextsep}{6pt plus 2pt minus 2pt}      % inline [h] floats
\setlength{\floatsep}{6pt plus 2pt minus 2pt}       % float <-> float (single-column)

\setlength{\dbltextfloatsep}{6pt plus 2pt minus 2pt} % text <-> double-column floats (figure*)
\setlength{\dblfloatsep}{6pt plus 2pt minus 2pt}     % float <-> float (double-column)

% --- caption spacing for tables only ---
\captionsetup[table]{aboveskip=6pt, belowskip=0pt}


\newenvironment{indentedquote}[1]%
{\list{}{\leftmargin=#1\rightmargin=#1}\item[]}%
{\endlist}


\newcommand\wordcount{\input{|"texcount -inc -sum -0 -utf8 -ch -template={SUM} \currfilepath"}}


%%%%%%%%%%%%%%%%%%%%%%%%%%%%%%%%%%%%%%%%%%%%%%%%%%%%%%%%%%%%%%%%%%%%%%%%%%%%%%%%

% title, author, institution

\title{Reliability-Aware Hierarchical Learning for Chagas Detection from Electrocardiogram under Expert Label Scarcity}
\author{Hao Wen\textsuperscript{1},
Jingsu Kang\textsuperscript{2} \\ \ \\
\textsuperscript{1}College of Science, China Agricultural University, Beijing, China\\
\textsuperscript{2}Tianjin Medical University, Tianjin, China
}

\begin{document}
\maketitle


%%%%%%%%%%%%%%%%%%%%%%%%%%%%%%%%%%%%%%%%%%%%%%%%%%%%%%%%%%%%%%%%%%%%%%%%%%%%%%%%

% abstract
\begin{abstract}

% almost finished

Aim: We present a reliability-aware hierarchical learning framework for ECG-based Chagas cardiomyopathy screening in the George B. Moody PhysioNet Challenge 2025 by Team Revenger, aiming to maximize positive case retrieval under prevalence constraints.

Methods: The 12-lead ECGs were resampled to 400 Hz, bandpass filtered (0.5–45 Hz), and z-score normalized. We used a ResNet model integrated with squeeze-and-excitation (SE) modules for binary classification. To address severe class imbalance and the scarcity of expert-confirmed labels, we applied stratified upsampling and reliability-weighted label smoothing to prioritize expert-confirmed positives over self-reported ones. Model training used an asymmetric loss to further penalize false negatives and was optimized with AdamW and a OneCycle learning rate scheduler. Model selection was based on the Challenge score from an internal hold-out subset.

Results: On the official hidden test set, our method received a Challenge score of 0.163, ranking 32nd of 40 eligible teams.

Conclusion: The proposed method demonstrates effective performance for ECG-based Chagas screening, highlighting its potential for improving detection accuracy and reliability in resource-limited scenarios.

\end{abstract}


% main body of the paper

\section{Introduction}
\label{sec:intro}

% almost finished

Addressing underdiagnosis of Chagas disease through scalable ECG-based screening is the focus of the 2025 George B. Moody PhysioNet Challenge \cite{goldberger2000physionet, 2025ChallengeCinC, 2025ChallengePreprint}. Enabled by aggregated multi-cohort ECG datasets \cite{ribeiro2020automatic, Cardoso_2016_Sami-Trop, wagner2020ptb_xl, Nunes_2021_Incidence, Filho_2020_Prognostic}, the Challenge frames a multi-source learning setting with heterogeneous label reliability and severe class imbalance.

In this work, we propose a reliability-aware hierarchical framework that prioritizes expert-confirmed labels and mitigates severe class imbalance within a deep ECG model, with optimization aligned to prevalence-constrained sensitivity objectives.


\section{Methods}
\label{sec:methods}

% almost finished

\subsection{Datasets and Preprocessing}
\label{subsec:data}

% almost finished

We used three ECG datasets for model training, with substantial differences in sample size, Chagas prevalence, and label provenance as summarized in Table~\ref{tab:dataset_stats}.

\setlength{\tabcolsep}{4pt} % Default value: 6pt
\begin{table}[!htp]
\centering
\begin{tabular}{rlll}
\toprule
Dataset & \multicolumn{1}{c}{Size} & \multicolumn{1}{c}{Chagas rate} & Label provenance \\
\midrule
SaMi-Trop   & 1\,631     & 100 \%    & expert-confirmed \\
CODE-15\%   & 345\,779   & 1.795 \%  & self-reported \\
PTB-XL      & 21\,799    & 0 \%      & N/A \\
\bottomrule
\end{tabular}
\caption{Dataset statistics and label provenance. Chagas rate is the proportion of recordings labeled positive in each dataset. N/A indicates that confirmed Chagas cases are not expected (non-endemic population).}
\label{tab:dataset_stats}
\end{table}
\setlength{\tabcolsep}{6pt} % Default value: 6pt

All ECGs were uniformly resampled to 400 Hz, bandpass filtered (0.5–45 Hz), and z-score normalized to zero mean and unit variance computed as in Eq.~\ref{eq:z_score}:
\begin{equation}
\label{eq:z_score}
\tilde{\mathbf{x}} = \frac{\mathbf{x} - \boldsymbol{\mu}_{\mathbf{x}}}{\boldsymbol{\sigma}_{\mathbf{x}}}
\end{equation}
where $\mathbf{x}$ is the original ECG signal, $\boldsymbol{\mu}_{\mathbf{x}}$ and $\boldsymbol{\sigma}_{\mathbf{x}}$ are the mean and standard deviation of $\mathbf{x}$, respectively. We excluded ECGs shorter than 1200 samples to ensure inputs contain enough cardiac cycles for stable model analysis.


\begin{figure*}[!t]
\centering
\adjustbox{width=0.73\linewidth}{
\begin{tikzpicture}[
    % font=\small,
    line width=0.75pt,
    >={Latex[length=3mm]},
    node distance=6.4mm,
    every node/.style={inner sep=2pt},
    stem/.style={draw, rounded corners=2pt, minimum width=30mm, minimum height=9.5mm,
                 align=center, fill=blue!10},
    block/.style={draw, rounded corners=2pt, minimum width=30mm, minimum height=12mm,
                  align=center, fill=orange!12},
    se/.style={draw, rounded corners=2pt, minimum width=30mm, minimum height=9.5mm,
               align=center, fill=green!12},
    pool/.style={draw, rounded corners=2pt, minimum width=30mm, minimum height=9.5mm,
                 align=center, fill=cyan!10},
    head/.style={draw, rounded corners=2pt, minimum width=30mm, minimum height=9.5mm,
                 align=center, fill=purple!10},
    output/.style={draw, rounded corners=2pt, minimum width=30mm, minimum height=9.5mm,
                align=center, fill=red!10},
    arrow/.style={-{Latex[length=3mm]}, thick},
    darr/.style={-{Latex[length=2.4mm]}, thick},
    dashedarrow/.style={-{Latex[length=2.6mm]}, thick, dashed},
    ann/.style={inner sep=1pt, align=center},
    layer/.style={draw, rounded corners=2pt, minimum width=19mm, minimum height=6.2mm,
                  align=center, fill=blue!6},
    bnact/.style={draw, rounded corners=2pt, minimum width=19mm, minimum height=6.2mm,
                  align=center, fill=violet!10},
    act/.style={draw, rounded corners=2pt, minimum width=19mm, minimum height=6.2mm,
                  align=center, fill=red!10},
    drop/.style={draw, rounded corners=2pt, minimum width=19mm, minimum height=6.2mm,
                  align=center, fill=gray!15, dashed},
    sum/.style={draw, circle, inner sep=1.6pt, fill=white},
    % brace/.style={decorate, decoration={brace, amplitude=6pt}},
    detailbox/.style={draw, rounded corners=3pt, inner sep=6pt, fill=orange!3}
]

% ================= Left (overall vertical path) =================
\node[stem] (input) {Input \\ (12$\times$L)};
\node[stem, below=of input] (stem) {Stem \\ Conv ks=15, 64 Ch \\ BN + ReLU};

% Compressed single bottleneck block (represents 4 stacked blocks)
% \node[block, below=of stem] (bottleneck) {$4 \times$ Bottleneck Block \\ ks: 1--15--1, stride: 1--4--1 \\ Ch: $512 \to 768 \to 1024 \to 1280$};

\node[block, below=of stem, minimum height=0pt] (bottleneck) {%
\begin{tikzpicture}[baseline=(ch.base), node distance=1mm]
\node[inner sep=2pt] (core)
{Bottleneck Block \\ ks: 1--15--1, stride: 1--4--1};
\draw[rounded corners=1pt] (core.north west) rectangle (core.south east);
\node[minimum width=2mm, left=of core] {\larger[3] $\mathbf{4\times}$};
\node[below left=1mm and -20mm of core.south, inner sep=2pt] (ch) {Ch: $512 \to 768 \to 1024 \to 1280$};
\end{tikzpicture}
};

% Global SE
\node[se, below=of bottleneck] (gse) {Global SE \\ reduction=8};

% Global pooling
\node[pool, below=of gse] (gap) {Global MaxPool};

% MLP Head
\node[head, below=of gap] (mlp) {MLP Head \\ (1280 $\rightarrow$ C=2)};

% Output
\node[output, below=of mlp] (out) {Output Logits};

% Arrows (overall path)
\draw[arrow] (input) -- (stem);
\draw[arrow] (stem) -- (bottleneck);
\draw[arrow] (bottleneck) -- (gse);
\draw[arrow] (gse) -- (gap);
\draw[arrow] (gap) -- (mlp);
\draw[arrow] (mlp) -- (out);

% ================= Right (expanded single bottleneck internals) =================
% Anchor for expansion (to the right of overall bottleneck block)
\coordinate (expandAnchor) at ($(input.east)+(50mm,0.2)$);

\node[layer, below=-0.5mm of expandAnchor] (c1) {Conv ks=1};
\node[below right=-1mm and 8mm of c1.south] {reduce};
\node[bnact, below=5mm of c1] (bn1) {BN};
\node[act, below=5mm of bn1] (a1) {ReLU};

\node[layer, below=5mm of a1] (c2) {Conv ks=15\\(stride 4)};
\node[bnact, below=5mm of c2] (bn2) {BN};
\node[act, below=5mm of bn2] (a2) {ReLU};
% \node[drop, below=5mm of a2] (drp) {Dropout 0.2};

\node[layer, below=5mm of a2] (c3) {Conv ks=1};
\node[below right=-1mm and 8mm of c3.south] {expand};
\node[bnact, below=5mm of c3] (bn3) {BN};

\node[sum, below=5mm of bn3] (add) {+};
\node[act, below=5mm of add] (outact) {ReLU};

% Projection/skip branch (since stride=4 or channel change)
\node[layer, right=10mm of c2, yshift=-7mm, minimum width=19mm] (proj) {Proj Conv \\ (stride 4)};
\node[bnact, below=5mm of proj, minimum width=19mm] (projbn) {BN};

% Connections internal
\draw[darr] (c1) -- (bn1);
\draw[darr] (bn1) -- (a1);
\draw[darr] (a1) -- (c2);
\draw[darr] (c2) -- (bn2);
\draw[darr] (bn2) -- (a2);
% \draw[darr] (a2) -- (drp);
\draw[darr] (a2) -- (c3);
\draw[darr] (c3) -- (bn3);
\draw[darr] (bn3) -- (add);
\draw[darr] (add) -- (outact);

% Skip path
\draw[darr] (c1.east) -| ($(proj.north)+(0,0mm)$);
\draw[darr] (proj) -- (projbn);
\draw[darr] (projbn.south) |- (add.east);

\begin{scope}[on background layer]
\node[detailbox, fit=(c1) (bn1) (a1) (c2) (bn2) (a2) (c3) (bn3) (proj) (projbn) (add) (outact)] (detail) {};
\end{scope}

\draw[dashed, line width=0.9pt] (bottleneck.north east) -- (detail.north west);
\draw[dashed, line width=0.9pt] (bottleneck.south east) -- (detail.south west);

\end{tikzpicture}

}
\caption{Model architecture. Left: overall network: a stem (Conv1d, kernel size 15, 64 output channels, Batch Normalization, ReLU) followed by four bottleneck residual blocks, a global SE module, global max pooling, and an MLP head producing $C = 2$ logits. Channel widths shown ($512 \to 768 \to 1024 \to 1280$) are the expanded channels.
Right: internal bottleneck structure. The middle convolution of kernel size 15 uses a stride of 4 for temporal downsampling; kernel size 1 convolutions reduce and then expand channels, and a projection convolution (kernel size 1, stride 4) aligns resolution and width for the residual path. For clarity, dropout layers present in the implementation are omitted. Abbreviations: ks kernel size; Ch channels; BN Batch Normalization; SE squeeze-and-excitation; MLP multi-layer perceptron.}
\label{fig:model}
\end{figure*}


\subsection{Reliability-Aware Hierarchical Supervision}
\label{subsec:reliability}

% almost finished

We introduce a hierarchical supervision scheme that encodes source reliability through stratified label smoothing and adaptive upsampling. Three reliability levels are defined: (1) expert-confirmed positives (SaMi-Trop, maximal trust), (2) self-reported samples (CODE-15\%, both positives and negatives, higher uncertainty), and (3) non-endemic negatives (PTB-XL, very low true prevalence but still mildly regularized). Given a one-hot label $\mathbf{y}$ ($[0, 1]$ for positives, $[1, 0]$ for negatives) and number of classes $C = 2,$ the smoothed target is computed as in Eq.~\ref{eq:label_smoothing}:
\begin{equation}
\label{eq:label_smoothing}
\tilde{\mathbf{y}} = (1 - \varepsilon) \cdot {\mathbf{y}} + \frac{\varepsilon}{C} \cdot {\mathbf{1}},
\end{equation}
where $\varepsilon$ is the smoothing factor which depends on the reliability level: 0.0 (SaMi positives), 0.6 (CODE-15\% positives \& negatives), 0.2 (PTB-XL negatives). This attenuates overconfident gradients for noisier or potentially misreported labels while preserving sharp supervision on expert-confirmed cases.

Severe class imbalance was mitigated by upsampling positives during training: positive samples from CODE-15\% were upsampled by a factor of 3, and those from SaMi-Trop by 12. No upsampling was applied to PTB-XL, which contains no positives. We chose these factors after reviewing hidden validation scores from multiple submissions, as shown in Table~\ref{tab:upsampling_schemes}.
\begin{table}[!htbp]
\centering
\begin{tabular}{ccc}
\toprule
\multicolumn{2}{c}{Upsample factor} & \multirow{2}{*}{Challenge score} \\ \cmidrule(lr){1-2}
CODE-15\% & SaMi-Trop               &                                  \\ \midrule
-         & -                       & 0.239$^\dagger$                        \\
3         & 7                       & 0.212                            \\
3         & 12                      & \textbf{0.245}                   \\
10        & 120                     & 0.210                            \\
6         & 36                      & 0.221                            \\
\bottomrule
\end{tabular}
\caption{Representative upsampling schemes and corresponding Challenge scores on the hidden validation set. The model and training strategies used were the same. “-” indicates no upsampling.\\
$^\dagger$ obtained during the unofficial phase.
}
\label{tab:upsampling_schemes}
\end{table}


Smoothed labels and upsampling strategies for each dataset are summarized in Table~\ref{tab:augmentation}.


% \begin{table}[!htp]
\begin{table}[!h]
\centering
\begin{tabular}{lcccc}
\toprule
\multirow{2}{*}{Dataset} & Upsampling & \multicolumn{2}{c}{Smoothed labels} \\ \cmidrule(lr){3-4}
& factor & negative & positive \\
\midrule
SaMi-Trop & $12$ & N/A          & $[0,\,1]$     \\
CODE-15\% & $3$ & $[0.7,\,0.3]$ & $[0.3,\,0.7]$ \\
PTB-XL    & $1$ & $[0.9,\,0.1]$ & N/A           \\
\bottomrule
\end{tabular}
\caption{Smoothed labels (computed from Eq.~\ref{eq:label_smoothing}) and upsampling strategies for each dataset.}
\label{tab:augmentation}
\end{table}


\subsection{Model Architecture}
\label{subsec:model}

% almost finished

We build upon the 1D ResNet ECG classifier of Ribeiro et al.~\cite{ribeiro2020automatic} and introduce three modifications.

\textbf{(1) Bottleneck residual blocks.} We replace basic ResNet blocks with bottleneck blocks of kernel sizes $1$–$15$–$1$ (pointwise--temporal--pointwise). The middle temporal convolution applies a stride of 4 for downsampling; the two convolutions (kernel size 1) first reduce the number of channels and then expand them with an expansion factor of 4. A projection convolution (kernel size 1, stride 4) is used in the residual branch whenever temporal resolution or channel width changes. Across the four blocks, the reduced (bottleneck) channel widths are $128 \to 192 \to 256 \to 320$, yielding expanded output widths $512 \to 768 \to 1024 \to 1280$.

\textbf{(2) Global squeeze-and-excitation (SE).} After the final bottleneck block, a single global SE module (reduction ratio 8) \cite{hu2018senet} performs temporal average pooling to a channel descriptor, applies a two-layer bottleneck multi-layer perceptron (MLP) $1280 \to 160 \to 1280$ with ReLU and sigmoid gating, and rescales the feature map channel-wise.

\textbf{(3) Global pooling head for variable input length.} Instead of flattening a fixed-length feature map as in the original baseline, we apply global max pooling over the remaining temporal dimension, yielding a 1280-dimensional vector irrespective of input length $L$. This vector is fed to a lightweight two-layer classification MLP: a hidden fully connected layer ($1280 \to 1024$) with non-linear activation and dropout (rate 0.2), followed by a final linear layer ($1024 \to 2$) producing class logits.

\textbf{(4) Stem and regularization.} A stem Conv1d (kernel size 15, stride 1, 64 channels) with BatchNorm and ReLU precedes the bottleneck stack. Within each bottleneck block, we apply BatchNorm+ReLU after the first two convolutions and dropout (rate 0.2) after each of those activations. All convolutions use “same” padding to preserve temporal length before downsampling operations.

% Overall, the changes primarily enable variable-length inference and add global channel reweighting with minimal structural overhead.
The overall model architecture is illustrated in Fig.~\ref{fig:model}.


\subsection{Training and Implementation Setups}
\label{subsec:train}

% almost finished

We employed asymmetric loss (ASL) \cite{ridnik2021asymmetric_loss} to complement the reliability-aware label smoothing strategy, jointly addressing the challenges of severe class imbalance. Let ${\mathbf{z}} = (z_0, z_1)$ denote the logits and $p=\operatorname{softmax}({\mathbf{z}})_1$ the predicted probability of the positive class. The ASL is defined in Eq.~\ref{eq:asl} with separate focusing parameters for positives and negatives and a clipped negative probability term:
\begin{equation}
\label{eq:asl}
\begin{multlined}
L = -y \cdot (1-p)^{\gamma_{+}} \log(p) \\
\phantom{L = } - (1-y) \cdot (p_m)^{\gamma_{-}} \log(1-p_m),
\end{multlined}
\end{equation}
where $y$ is the (smoothed) positive-class target probability, $p_m = \max(p - m, 0),$ $(\gamma_{+},\gamma_{-})=(1,4)$ and margin $m=0.05$. We train for 30 epochs with batch size 128 using the AdamW optimizer (initial learning rate $1\times10^{-4}$, peak $6\times10^{-4}$ under a OneCycle scheduler, weight decay $1\times10^{-2}$). Early stopping (patience 10 epochs, monitored on a fixed 20\% internal hold-out subset) selects the final model via the Challenge metric. Each training segment is a uniform random crop (or center padding if shorter) of length 4096 samples. The full implementation, including model construction, data pipeline, and optimization utilities, is based on the \texttt{torch-ECG} framework \cite{torch_ecg_paper}.


\section{Results}
\label{sec:results}

% almost finished

The Challenge score of our team ``Revenger'' on the hidden test set was 0.163, ranking 32nd among 40 eligible teams. This score and ranking, along with extra scores on the internal hold-out of the public training data, and on the hidden validation set, are summarized in Table~\ref{tab:scores}.

\begin{table}[!htp]
\centering
% \begin{tabular}{r|r|r|r}
Training          & Validation & Test & Ranking   \\ \hline
$0.451 \pm 0.005$ & 0.245      & 0.163  & 32~/~40 \\ \hline
\end{tabular}

\begin{tabular}{r|r|r|r}
Training          & Validation & \textbf{Test}  & \textbf{Ranking} \\ \hline
$0.451 \pm 0.005$ & 0.245      & \textbf{0.163} & \textbf{32~/~40} \\ \hline
\end{tabular}

\caption{Challenge scores for our submitted entries (team ``Revenger''). Training: internal hold-out mean $\pm$ std over repeated runs. Validation: best among 10 validation submissions. Test: the unique test submission. Ranking: position on the hidden test leaderboard.}
\label{tab:scores}
\end{table}


\section{Discussion and Conclusions}
\label{sec:discu}

% almost finished.

The Challenge scores presented in Table~\ref{tab:scores} indicate that our proposed method is effective for Chagas screening from ECGs, albeit with substantial room for improvement. The result demonstrates our model's ability to learn diagnostically relevant features from ECGs for this task under scarce and noisy supervision. This is achieved through reliability-aware label smoothing, which incorporates both label provenance and reliability instead of treating all positive labels uniformly. Together with the asymmetric loss and strategic upsampling, these results indicate that explicitly modeling label reliability helps stabilize the learning process more effectively than introducing additional architectural complexity. Overall, our approach offers a scalable and resource-efficient solution and aligns well with the Challenge's objective of identifying high-risk individuals under limited serological testing capacity.

% However, the performance gap between our internal hold-out subset and the hidden test set suggests two primary limitations. First, the reliability weights (smoothing factors) were pre-defined based on label provenance rather than learned from data. This static assignment cannot capture the inherent heterogeneity of label quality within each source. Second, positive upsampling was applied uniformly at the dataset level using fixed factors. This strategy overlooks variations in individual sample difficulty and does not adapt to the model's evolving confidence during training. Furthermore, we did not make use of the additional labels available in some datasets, such as arrhythmia labels, to design and implement auxiliary learning tasks that could have enhanced the performance of the main task, Chagas screening.

However, the performance gap between our internal hold-out subset and the hidden test set indicates limitations in our method's generalization capability. This is also reflected in the absolute performance, where our Challenge score is substantially behind those of the top-performing teams (e.g., 0.323, 0.283, 0.280). The leading approaches generally leveraged some form of large-scale representation learning, such as pre-trained Vision Transformer foundation models or self-supervised learning on extensive ECG datasets. These strategies focus on building robust, general-purpose feature extractors. In contrast, our work prioritized a distinct niche by designing a resource-efficient pipeline that explicitly addresses the challenges of label noise and scarcity through reliability-aware smoothing and asymmetric loss, without relying on massive pretraining. While this choice enhances practicality and stability under noisy supervision, it appears that the representational capacity of our directly-trained model is ultimately lower, limiting its ability to generalize as effectively as the foundation-model-based approaches. Furthermore, our static assignment of reliability weights, which cannot adapt to instance-specific label quality, represents another limitation compared to more dynamic or learned weighting schemes.

% Future research directions will primarily focus on developing a self-adaptive supervision framework. This includes dynamic weighting schemes for learning instance-specific reliability scores, moving beyond static smoothing factors, and adaptive sampling strategies that respond to the model's evolving confidence during training, offering a promising alternative to fixed upsampling factors. Data augmentation techniques, such as CutMix \cite{yun2019cutmix} and SMOTE \cite{Chawla_2002_SMOTE}, could be applied to further expand and diversify the positive samples, thereby enhancing the model's robustness against overfitting and improving generalization to underrepresented patterns. Additionally, a multi-task learning framework leveraging auxiliary arrhythmia labels could enhance feature representation and improve generalization for the primary Chagas screening task. Exploring self-supervised pre-training on large-scale unlabeled ECG data also represents a promising direction to learn more transferable representations before fine-tuning on the target task.

Building upon these insights, future research will focus on bridging the generalization gap while retaining our focus on learning under weak supervision. A primary direction is the self-adaptive supervision framework. This includes dynamic weighting schemes, moving beyond our current static factors, and adaptive sampling strategies that respond to the model's evolving confidence during training. To further improve robustness and handle class imbalance, advanced data augmentation techniques such as CutMix \cite{yun2019cutmix} and SMOTE \cite{Chawla_2002_SMOTE} will be explored to diversify the limited positive samples, with a particular focus on ECG-specific transformations like lead-wise masking. Furthermore, inspired by the success of representation learning, we plan to explore self-supervised pre-training on large-scale unlabeled ECG data as a pivotal step to learn more transferable representations before fine-tuning. This direction, while computationally more demanding, addresses a key limitation identified in our current work. Additionally, a multi-task learning framework leveraging auxiliary arrhythmia labels could be integrated to impose clinically meaningful constraints and enhance the feature representation for the primary Chagas screening task.


%%%%%%%%%%%%%%%%%%%%%%%%%%%%%%%%%%%%%%%%%%%%%%%%%%%%%%%%%%%%%%%%%%%%%%%%%%%%%%%%

% acknowledgments

% \section*{Acknowledgments}

% finished

The authors thank Professor Deren Han from Beihang University for helping accomplish this work.


%%%%%%%%%%%%%%%%%%%%%%%%%%%%%%%%%%%%%%%%%%%%%%%%%%%%%%%%%%%%%%%%%%%%%%%%%%%%%%%%

\bibliographystyle{cinc}
\bibliography{references}
% \bibliography{references-reduced}

\begin{correspondence}
Hao Wen\\
No. 17, Qinghua East Road, Haidian District, Beijing, China\\
wenh06@cau.edu.cn,wenh06@gmail.com
\end{correspondence}

\balance

\end{document}



%%%%%%%%%%%%%%%%%%%%%%%%%%%%%%%%%%%%%%%%%%%%%%%%%%
% conference poster
%%%%%%%%%%%%%%%%%%%%%%%%%%%%%%%%%%%%%%%%%%%%%%%%%%

% %%%%%%%%%%%%%%%%%%%%%%%%%%%%%%%%%%%%%%%%%
% NIWeek 2014 Poster by T. Reveyrand
% www.microwave.fr
% http://www.microwave.fr/LaTeX.html
% ---------------------------------------
%
% Original template created by:
% Brian Amberg (baposter@brian-amberg.de)
%
% This template has been downloaded from:
% http://www.LaTeXTemplates.com
%
% License:
% CC BY-NC-SA 3.0 (http://creativecommons.org/licenses/by-nc-sa/3.0/)
%
%%%%%%%%%%%%%%%%%%%%%%%%%%%%%%%%%%%%%%%%%

%----------------------------------------------------------------------------------------
%	PACKAGES AND OTHER DOCUMENT CONFIGURATIONS
%----------------------------------------------------------------------------------------

\documentclass[a0paper,portrait]{baposter}

\usepackage[T1]{fontenc}
\usepackage{mathtools, graphicx, xurl, hyperref}
% \usepackage{tikz-cd}
\usepackage{tikz}
\usetikzlibrary{shapes, arrows.meta, decorations.pathmorphing, decorations.pathreplacing, backgrounds, positioning, fit, calc, shadows, shapes.misc}
% \tikzexternalize[prefix=tikz/,optimize command away=\includepdf]
\usepackage{boldline, multirow}
\usepackage{booktabs} % Allows the use of \toprule, \midrule and \bottomrule in tables
\usepackage{colortbl}
\usepackage{subcaption}
\usepackage{relsize}
\usepackage{fontawesome5}
\usepackage[ruled, vlined]{algorithm2e}

\usepackage{tabularx} % for 'X' col. type and 'tabularx' environment
\usepackage{ragged2e} % for '\RaggedRight' macro
\newcolumntype{L}{>{\RaggedRight}X} % disable full justification

\usepackage{adjustbox}

% \tikzstyle{smallcircleblock} = [circle, draw, text width = 0.7em, text centered]

% \tikzstyle{wideblock} = [rectangle, draw, text width = 6.5em, text centered, rounded corners, inner sep = 7pt, minimum height = 1.0em]

\newcommand\ImageNode[4][]{
  \node[#1] (#2) {\includegraphics[#3]{#4}};
}
\definecolor{arrowblue}{RGB}{98,145,224}


\newcolumntype{g}{>{\columncolor[gray]{0.8}}l}


\newif\ifcoloredtext
\coloredtextfalse

\newif\ifboxednn
\boxednntrue


% --- unify spacing between floats and text (single-/double-column) ---
\setlength{\textfloatsep}{6pt plus 2pt minus 2pt}   % text <-> single-column floats (table/figure)
\setlength{\intextsep}{6pt plus 2pt minus 2pt}      % inline [h] floats
\setlength{\floatsep}{6pt plus 2pt minus 2pt}       % float <-> float (single-column)

\setlength{\dbltextfloatsep}{6pt plus 2pt minus 2pt} % text <-> double-column floats (figure*)
\setlength{\dblfloatsep}{6pt plus 2pt minus 2pt}     % float <-> float (double-column)

% --- caption spacing for tables only ---
\captionsetup[table]{aboveskip=6pt, belowskip=0pt}


\newenvironment{indentedquote}[1]%
{\list{}{\leftmargin=#1\rightmargin=#1}\item[]}%
{\endlist}


\newcommand\wordcount{\input{|"texcount -inc -sum -0 -utf8 -ch -template={SUM} \currfilepath"}}


\usepackage{textcomp}
\usepackage{eso-pic}
\usepackage{enumitem}
\usepackage{hyperref}
\usepackage{fix-cm}

\setlist[itemize]{
  itemsep=1.5pt,
  topsep=3pt,
  parsep=0pt,
  partopsep=0pt,
  leftmargin=1em,
  labelsep=0.4em
}

\setlength{\abovedisplayskip}{1.5pt plus 1pt minus 1pt}
\setlength{\belowdisplayskip}{1.5pt plus 1pt minus 1pt}

\graphicspath{{images/}} % Directory in which figures are stored

\definecolor{bordercol}{RGB}{40,40,40} % Border color of content boxes
\definecolor{headercol1}{RGB}{186,215,230} % Background color for the header in the content boxes (left side)
\definecolor{headercol2}{RGB}{120,120,120} % Background color for the header in the content boxes (right side)
\definecolor{headerfontcol}{RGB}{0,0,0} % Text color for the header text in the content boxes
\definecolor{boxcolor}{RGB}{210,235,250} % Background color for the content in the content boxes

\definecolor{jdhblue}{RGB}{2,93,186}

\definecolor{zkblue}{RGB}{88,135,175}
\definecolor{zkbackground}{RGB}{230,232,234}


% \usetikzlibrary{shapes, arrows, external, decorations.pathmorphing, backgrounds, positioning, fit, petri, calc, hobby, cd}


\begin{document}

% \background{ % Set the background to an image (background.pdf)
% \begin{tikzpicture}[remember picture,overlay]
% \draw (current page.north west)+(-2em,2em) node[anchor=north west]
% {\includegraphics[height=1.1\textheight]{images/baposter_background.pdf}};
% \end{tikzpicture}
% }

% add conference banner
\AddToShipoutPictureFG*{%
  \AtPageLowerLeft{%
    \hspace{0.77em}
    \raisebox{0.2ex}{\includegraphics[width=0.147\textwidth]{cinc2025-banner.png}}
  }
}

\begin{poster}{
grid=false,
borderColor=bordercol, % Border color of content boxes
headerColorOne=headercol1, % Background color for the header in the content boxes (left side)
headerColorTwo=headercol2, % Background color for the header in the content boxes (right side)
headerFontColor=headerfontcol, % Text color for the header text in the content boxes
% boxColorOne=boxcolor, % Background color for the content in the content boxes
boxColorOne=zkbackground,
headershape=roundedright, % Specify the rounded corner in the content box headers
headerfont=\Large\sf\bf, % Font modifiers for the text in the content box headers
textborder=rectangle,
% background=user,
background=none,
headerborder=open, % Change to closed for a line under the content box headers
boxshade=plain
}
%
%----------------------------------------------------------------------------------------
%	TITLE AND AUTHOR NAME
%----------------------------------------------------------------------------------------
%
{\includegraphics[scale=0.385]{logo_cau_science.png}} % University/lab logo
{
{\bf \fontsize{17pt}{17pt} \selectfont Reliability-Aware Hierarchical Learning for Chagas Detection from ECG under Expert Label Scarcity}
} % Poster title
{\vspace{0.3em} \smaller Hao WEN$^1$, Jingsu KANG$^2$  \\  % Author names

$^1${\it College of Science, China Agricultural University}\qquad
$^2${\it Tianjin Medical University} \\
\vspace{0.2cm}
{\Large \bf{The George B. Moody PhysioNet Challenge 2025}, ~~~\bf{Team Revenger}}
}
{\includegraphics[scale=0.097]{logo_tmu.jpeg}}


\newif\ifcoloredtext
\coloredtexttrue
\newif\ifboxednn
\boxednntrue


%----------------------------------------------------------------------------------------
%	INTRODUCTION & CHALLENGE CONTEXT
%----------------------------------------------------------------------------------------

\headerbox{Framework Summary \& Key Components}{name=introduction, column=0, row=0, span=3}{

\begin{itemize}
\item We propose a \textbf{reliability-aware hierarchical framework} for ECG-based Chagas screening under scarce and noisy labels. The framework comprises three coordinated components: (1) provenance‑stratified label smoothing, (2) moderated source-specific upsampling, and (3) a lightweight variable‑length bottleneck ResNet classifier.
\item The supervision design assigns sharp targets to expert‑confirmed positives, softened targets to self‑reported samples, and mild smoothing to non‑endemic negatives, while distinct upsampling factors amplify scarce reliable positives without collapsing cohort diversity.
\item The encoder stacks 4 bottleneck residual blocks with a global squeeze‑and‑excitation module and global max pooling to handle variable-length inputs; training couples reliability‑coded targets with an asymmetric loss to concentrate learning on hard positives amid severe class imbalance.
\end{itemize}
}

%----------------------------------------------------------------------------------------
%	DATA
%----------------------------------------------------------------------------------------

\headerbox{Dataset Preparation}{name=data, column=0, row=1, span=1, below=introduction}{

\begin{itemize}
\item \textbf{Sampling}: All 12-lead ECGs uniformly resampled to 400 Hz.
\item \textbf{Filtering}: Butterworth bandpass 0.5–45 Hz to suppress baseline wander + high-frequency noise.
\item \textbf{Normalization}: Per-record z-score
\begingroup
\setlength{\abovedisplayskip}{2pt}  % 上间距
\setlength{\belowdisplayskip}{2pt}  % 下间距
\[
\tilde{\mathbf{x}} = \frac{\mathbf{x} - \boldsymbol{\mu}_{\mathbf{x}}}{\boldsymbol{\sigma}_{\mathbf{x}}}
\]
\endgroup
to zero mean, unit variance.
\item \textbf{Exclusion criteria}: Discard recordings $<1{,}200$ samples ($<3$ s @ 400 Hz) to ensure enough cardiac cycles.
\end{itemize}

}

%----------------------------------------------------------------------------------------
%	RELIABILITY-AWARE SUPERVISION
%----------------------------------------------------------------------------------------

\headerbox{Reliability-Aware Supervision}{name=supervision, column=1, row=1, span=2, below=introduction}{

% Three ECG datasets with heterogeneous label reliability and Chagas prevalence:

Hierarchical supervision: source reliability levels (expert > self-report > non-endemic negatives) $\longrightarrow$ calibrated label smoothing ($\varepsilon$) + moderated upsampling.

\vspace{-1ex}
% {\smaller[1]
\begin{center}
\begin{tabular}{rlllllll}
\toprule
\multirow{2}{*}{Dataset} & \multirow{2}{*}{Size} & \multirow{2}{*}{Chagas \%} & \multirow{2}{*}{Provenance} & \multirow{2}{*}{Up $\times$} & \multirow{2}{*}{$\varepsilon$} & \multicolumn{2}{c}{Smoothed Targets} \\ \cmidrule(lr){7-8}
& & & & & & negative & positive \\
\midrule
SaMi-Trop   & 1,631     & 100 \%    & expert-confirmed  & $\times$ 12 & 0.0 & N/A           & $[0,\,1]$     \\
CODE-15\%   & 345,779   & 1.795 \%  & self-reported     & $\times$ 3  & 0.6 & $[0.7,\,0.3]$ & $[0.3,\,0.7]$ \\
PTB-XL      & 21,799    & 0 \%      & non-endemic (neg) & $\times$ 1  & 0.2 & $[0.9,\,0.1]$ & N/A           \\
\bottomrule
\end{tabular}
\end{center}
% }
\vspace{-1ex}

Label smoothing (binary, $C=2$): $\tilde{\mathbf{y}} = (1 - \varepsilon) \cdot {\mathbf{y}} + \frac{\varepsilon}{C} \cdot {\mathbf{1}}$.

\begin{itemize}
% \item Reliability levels: expert (sharp, high upsampling) > self-report (soft, moderate) > non-endemic negatives (mild).
\item Purpose: reduce overconfident gradients for noisier self-reports while preserving expert signal.
% \item Upsampling balance: (CODE-15\% $\times$3, SaMi-Trop $\times$12) maximized hidden score; aggressive (10,120) overfit / degraded generalization.
% \item PTB-XL adds negative morphology diversity; mild smoothing (0.2) regularizes the decision boundary.
% \item Interaction: provenance-aware smoothing complements asymmetric loss (ASL) by encoding prior trust.
\item Benefit: replaces ad-hoc class weights with structured priors; stabilizes early training variance.
% \item Limitation: $\varepsilon$ values are manually fixed—do not capture intra-source heterogeneity (future: learn instance reliability).
\item Complementarity: provenance-based smoothing reduces label noise impact; ASL (asymmetric loss) focuses capacity on hard positives:
\begingroup
\setlength{\abovedisplayskip}{2pt}  % 上间距
\setlength{\belowdisplayskip}{2pt}  % 下间距
\[
\operatorname{ASL} = -y \cdot (1-p)^{\gamma_{+}} \log(p) - (1-y) \cdot (p_m)^{\gamma_{-}} \log(1-p_m)
\]
\endgroup
where $p$ is the predicted probability of the positive class, $y$ is the (smoothed) positive-class target probability, $p_m = \max(p - m, 0),$ $(\gamma_{+},\gamma_{-})=(1,4)$ and margin $m=0.05$.
\end{itemize}

}

%----------------------------------------------------------------------------------------
%	MODEL ARCHITECTURE
%----------------------------------------------------------------------------------------

\headerbox{Model Architecture}{name=model, column=0, row=2, span=1, below=data}{

{\bf Enhanced 1D ResNet with SE Modules}

\vspace{-2.5ex}

\begin{center}
\resizebox{1.03\linewidth}{!}{%
\begin{tikzpicture}[
    % font=\small,
    line width=0.75pt,
    >={Latex[length=3mm]},
    node distance=6.4mm,
    every node/.style={inner sep=2pt},
    stem/.style={draw, rounded corners=2pt, minimum width=30mm, minimum height=10mm,
                 align=center, fill=blue!10},
    block/.style={draw, rounded corners=2pt, minimum width=30mm, minimum height=12mm,
                  align=center, fill=orange!12},
    se/.style={draw, rounded corners=2pt, minimum width=30mm, minimum height=10mm,
               align=center, fill=green!12},
    pool/.style={draw, rounded corners=2pt, minimum width=30mm, minimum height=10mm,
                 align=center, fill=cyan!10},
    head/.style={draw, rounded corners=2pt, minimum width=30mm, minimum height=10mm,
                 align=center, fill=purple!10},
    output/.style={draw, rounded corners=2pt, minimum width=24mm, minimum height=10mm,
                align=center, fill=red!10},
    arrow/.style={-{Latex[length=3mm]}, thick},
    darr/.style={-{Latex[length=2.4mm]}, thick},
    dashedarrow/.style={-{Latex[length=2.6mm]}, thick, dashed},
    ann/.style={inner sep=1pt, align=center},
    layer/.style={draw, rounded corners=2pt, minimum width=18mm, minimum height=6.2mm,
                  align=center, fill=blue!6},
    bnact/.style={draw, rounded corners=2pt, minimum width=18mm, minimum height=6.2mm,
                  align=center, fill=violet!10},
    act/.style={draw, rounded corners=2pt, minimum width=18mm, minimum height=6.2mm,
                  align=center, fill=red!10},
    drop/.style={draw, rounded corners=2pt, minimum width=18mm, minimum height=6.2mm,
                  align=center, fill=gray!15, dashed},
    sum/.style={draw, circle, inner sep=1.6pt, fill=white},
    % brace/.style={decorate, decoration={brace, amplitude=6pt}},
    detailbox/.style={draw, rounded corners=3pt, inner sep=6pt, fill=orange!3}
]

% ================= Left (overall vertical path) =================
\node[stem] (input) {Input \\ (12$\times$L)};
\node[stem, below=of input] (stem) {Stem \\ Conv ks=15 \\ BN + ReLU \\ 64 Ch};

% Compressed single bottleneck block (represents 4 stacked blocks)
\node[block, below=of stem] (bottleneck) {$4 \times$ Bottleneck Block \\ ks: 1--15--1 (stride=4) \\ Ch: $512 \to 768 \to 1024 \to 1280$};

% Global SE
\node[se, below=of bottleneck] (gse) {Global SE \\ reduction=8};

% Global pooling
\node[pool, below=of gse] (gap) {Global MaxPool};

% MLP Head
\node[head, below=of gap] (mlp) {MLP Head \\ (1280 $\rightarrow$ C=2)};

% Output
\node[output, below=of mlp] (out) {Output Logits};

% Arrows (overall path)
\draw[arrow] (input) -- (stem);
\draw[arrow] (stem) -- (bottleneck);
\draw[arrow] (bottleneck) -- (gse);
\draw[arrow] (gse) -- (gap);
\draw[arrow] (gap) -- (mlp);
\draw[arrow] (mlp) -- (out);

% ================= Right (expanded single bottleneck internals) =================
% Anchor for expansion (to the right of overall bottleneck block)
\coordinate (expandAnchor) at ($(input.east)+(45mm,0)$);

\node[layer, below=0mm of expandAnchor] (c1) {Conv 1$\times$1};
\node[below right=-1mm and 8mm of c1.south] {reduce};
\node[bnact, below=5mm of c1] (bn1) {BN};
\node[act, below=5mm of bn1] (a1) {ReLU};

\node[layer, below=5mm of a1] (c2) {Conv 15\\(stride 4)};
\node[bnact, below=5mm of c2] (bn2) {BN};
\node[act, below=5mm of bn2] (a2) {ReLU};
% \node[drop, below=5mm of a2] (drp) {Dropout 0.2};

\node[layer, below=5mm of a2] (c3) {Conv 1$\times$1};
\node[below right=-1mm and 8mm of c3.south] {expand};
\node[bnact, below=5mm of c3] (bn3) {BN};

\node[sum, below=5mm of bn3] (add) {+};
\node[act, below=5mm of add] (outact) {ReLU (output)};

% Projection/skip branch (since stride=4 or channel change)
\node[layer, right=10mm of c2, yshift=-7mm, minimum width=18mm] (proj) {Proj 1$\times$1 \\ (stride 4)};
\node[bnact, below=5mm of proj, minimum width=18mm] (projbn) {BN};

% Connections internal
\draw[darr] (c1) -- (bn1);
\draw[darr] (bn1) -- (a1);
\draw[darr] (a1) -- (c2);
\draw[darr] (c2) -- (bn2);
\draw[darr] (bn2) -- (a2);
% \draw[darr] (a2) -- (drp);
\draw[darr] (a2) -- (c3);
\draw[darr] (c3) -- (bn3);
\draw[darr] (bn3) -- (add);
\draw[darr] (add) -- (outact);

% Skip path
\draw[darr] (c1.east) -| ($(proj.north)+(0,0mm)$);
\draw[darr] (proj) -- (projbn);
\draw[darr] (projbn.south) |- (add.east);

\begin{scope}[on background layer]
\node[detailbox, fit=(c1) (bn1) (a1) (c2) (bn2) (a2) (c3) (bn3) (proj) (projbn) (add) (outact)] (detail) {};
\end{scope}

\draw[dashed, line width=0.9pt] (bottleneck.north east) -- (detail.north west);
\draw[dashed, line width=0.9pt] (bottleneck.south east) -- (detail.south west);

\end{tikzpicture}

}
\end{center}

\vspace{-0.5ex}

\begin{itemize}
\item {\bf Stem}: Conv1d (kernel 15, stride 1, 64 output channels) followed by BN and ReLU, providing initial local morphology modeling.
\item {\bf Bottleneck blocks}: 1-15-1 convolutions, expansion factor 4, stride 4 downsampling
\item {\bf Global SE}: Channel attention, reduction ratio 8 for temporal feature recalibration.
\item {\bf Variable-length inputs}: Global MaxPool enables inference on ECGs of any duration.
% \item {\bf Classification MLP head}: Hidden FC layer ($1280 \to 1024$) with non-linear activation and dropout (rate 0.2), followed by a final linear layer ($1024 \to 2$) producing class logits.
\end{itemize}

}

%----------------------------------------------------------------------------------------
%	UPSAMPLING STRATEGY SELECTION
%----------------------------------------------------------------------------------------

\headerbox{Upsample Strategy Selection}{name=upsampling, column=1, row=2, span=1, below=supervision}{
\small
% \textbf{Goal:} Amplify scarce reliable positives without collapsing inter‑cohort diversity.

\begin{center}
\begin{tabular}{ccc}
\toprule
\multicolumn{2}{c}{Upsample factor} & \multirow{2}{*}{Challenge score} \\ \cmidrule(lr){1-2}
CODE-15\% & SaMi-Trop               &                                  \\ \midrule
-         & -                       & 0.239$^\dagger$                        \\
3         & 7                       & 0.212                            \\
3         & 12                      & \textbf{0.245}                   \\
10        & 120                     & 0.210                            \\
6         & 36                      & 0.221                            \\
\bottomrule
\end{tabular}
\end{center}
\vspace{-1ex}
$^\dagger$ obtained during the unofficial phase.

\begin{itemize}
\item Moderate replication outperformed both no upsampling and extreme factors.
\item Extreme factors increased memorization and reduced hidden generalization.
% \item PTB-XL (negatives only) kept at ×1 to avoid further skew toward the negative class.
\end{itemize}

}

%----------------------------------------------------------------------------------------
%	TRAINING CONFIGURATION
%----------------------------------------------------------------------------------------

\headerbox{Training Configuration}{name=training, column=2, row=2, span=1, below=supervision}{

\begin{itemize}[leftmargin=8pt,itemsep=2pt,topsep=2pt]
\item \textbf{Optimizer \& scheduler}: AdamW (weight decay $1\times 10^{-2}$) with OneCycle LR over 30 epochs: base $1\times10^{-4} \rightarrow$ peak $6\times10^{-4}$ (early phase) $\rightarrow$ anneal to $\approx 1\times10^{-6}$.
% \item \textbf{Batch / epochs / early stop}: Batch size 128; max 30 epochs; early stopping (patience 10) on a fixed 20\% hold-out subset using the Challenge metric.
\item \textbf{Batch size}: 128.
\item \textbf{Epochs}: Max 30 epochs.
\item \textbf{Early stop}: Patience 10 epochs, evaluated on a fixed 20\% hold-out subset using the Challenge metric.
\item \textbf{Training segments}: Uniform random crop (or center padding if shorter) of 4{,}096 samples.
\item \textbf{Implementation}: PyTorch + \texttt{torch-ECG} framework

\end{itemize}
}

%----------------------------------------------------------------------------------------
%	RESULTS
%----------------------------------------------------------------------------------------

\headerbox{Results}{name=results, column=1, row=3, span=2, below=upsampling}{

% {\bf Challenge Performance \& Analysis}

\begin{center}
\begin{tabular}{r|r|r|r}
Internal hold-out (mean±std) & Hidden validation & Test & Ranking \\
\hline
$0.451 \pm 0.005$ & 0.245 & 0.163 & 32~/~40 \\
\hline
\end{tabular}
\end{center}
\vspace{-1ex}
The rankings are based on the hidden test set. Results on the public training data (more specifically, the hold-out subset) are also provided for reference.

}

%----------------------------------------------------------------------------------------
%	DISCUSSION & Limitations
%----------------------------------------------------------------------------------------

\headerbox{Discussion \& Limitations}{name=discussion, column=0, row=4, span=3, below=results}{

\begin{itemize}
\item Our reliability-aware hierarchical framework provides a lightweight, supervision-centric route to Chagas ECG screening under scarce positives and mixed label noise, without relying on heavy architectures or pretraining.
\item The generalization gap (internal hold-out vs hidden validation) suggests residual distribution shift (cohort composition/prevalence) and under-adaptation of our fixed supervision scheme.
\item Static supervision limits adaptation: (i) Fixed smoothing factors ($\varepsilon$) capture only coarse provenance, missing intra-source label quality variation; (ii) dataset-level upsampling with fixed multipliers ignores evolving model confidence and per-instance difficulty.
\item \textbf{Future directions}:
\begin{itemize}
\item Adaptive supervision: jointly learn instance-level reliability weights and adjust positive sampling (curriculum/difficulty- or confidence-aware, hard positive mining, etc.) instead of fixed smoothing factors $\varepsilon$ and static upsampling multipliers.
\item Augmentation \& synthesis: CutMix, mixup variants, temporal warping, and class-conditional synthesis to expand minority morphology diversity.
\item Multi-task \& pretraining: leverage arrhythmia / morphology labels; self-supervised ECG representation learning before supervised fine-tuning.
\end{itemize}
\end{itemize}

}

%----------------------------------------------------------------------------------------
%	FOOTER
%----------------------------------------------------------------------------------------

\headerbox{}{name=foottext, column=0, span=3, below=discussion, textborder=none, headerborder=none, boxheaderheight=0pt}{
\hfill \small \textit{Code, configs, etc. are available at https://github.com/wenh06/cinc2025, and is based on https://github.com/DeepPSP/torch\_ecg.}
}

\end{poster}

\end{document}



%%%%%%%%%%%%%%%%%%%%%%%%%%%%%%%%%%%%%%%%%%%%%%%%%%
% long paper
%%%%%%%%%%%%%%%%%%%%%%%%%%%%%%%%%%%%%%%%%%%%%%%%%%

% \input{paper_long}
